%======================================================================
\section{The ledger}
\label{sec:ledger}

This section introduces the paper's central computational object and
derives, from it alone, every pricing formula the rest of the paper uses.

\subsection{Definition and first properties}

\begin{definition}[Upper-share ledger]
\label{def:ledger}
For an admissible slice, define
\begin{equation}
G(z)\;=\;\int_{z}^{\infty}e^{\,x(t)}\,\rho(t)\,\dd t
\;=\;\E\big[\,Y\,\1\{Z>z\}\,\big] .
\label{eq:ledger}
\end{equation}
\end{definition}

$G(z)$ is the share of the (unit) mean carried by outcomes above the
$z$-th log-odds level --- hence the name.  Three properties are immediate
and get used constantly:
\begin{equation}
G(-\infty)=\E[Y]=1,
\qquad
G(+\infty)=0,
\qquad
G'(z)=-e^{\,x(z)}\rho(z)\;<\;0 .
\label{eq:ledgerprops}
\end{equation}
The ledger is a smooth, strictly decreasing curve from $1$ to $0$, whose
derivative is explicitly known at every node of whatever grid carries it.
That last fact looks small and is not: it is what will make interpolation
fourth-order accurate without any spline solve
(\cref{sec:computation}), and what makes the cancellation theorem of
\cref{sec:jacobian} one line.

For a strike $k$, define the \emph{strike root} $z_k$ and its rank $u_k$:
\begin{equation}
x(z_k)=k,
\qquad
u_k=\Lambda(z_k)=F_X(k) ,
\label{eq:strikeroot}
\end{equation}
unique because $x$ is a strictly increasing bijection of $\R$.

\subsection{Pricing identities}

\begin{proposition}[Ledger pricing]
\label{prop:pricing}
For every admissible slice and every $k$:
\begin{align}
c(k)&=G(z_k)-e^{k}\big(1-u_k\big),
\label{eq:call}\\
P(k)&=e^{k}u_k-\big(1-G(z_k)\big),
\label{eq:put}\\
c'(k)&=-e^{k}\big(1-u_k\big),
\label{eq:cprime}\\
c''(k)&=-e^{k}\big(1-u_k\big)+e^{k}\,\frac{\rho(z_k)}{x'(z_k)},
\label{eq:csecond}\\
f_X(k)&=e^{-k}\big(c''(k)-c'(k)\big)=\frac{\rho(z_k)}{x'(z_k)} .
\label{eq:bl_division}
\end{align}
\end{proposition}

\begin{proof}
\eqref{eq:call}: split the payoff on the event $\{Y>e^{k}\}=\{Z>z_k\}$
(the events coincide because $x$ is increasing),
\[
c(k)=\E\big[(Y-e^{k})\1\{Z>z_k\}\big]
=\E\big[Y\1\{Z>z_k\}\big]-e^{k}\,\Prob(Z>z_k)
=G(z_k)-e^{k}(1-u_k):
\]
the asset delivered on exercise, minus the strike bill paid on the same
event --- a call price is the difference of two ledgers, both monotone,
both already on the grid.  \eqref{eq:put} is the same split on the
complementary event, and subtracting gives parity
$c-P=G(z_k)+\big(1-G(z_k)\big)-e^{k}=1-e^{k}$, as it must.

\eqref{eq:cprime}: differentiate \eqref{eq:call} in $k$.  The chain rule
through the moving root contributes
$\big[G'(z_k)+e^{k}\rho(z_k)\big]\frac{\dd z_k}{\dd k}$; but
$G'(z_k)=-e^{x(z_k)}\rho(z_k)=-e^{k}\rho(z_k)$ by \eqref{eq:ledgerprops}
and \eqref{eq:strikeroot}, so the bracket vanishes identically, and only
the explicit $-e^{k}(1-u_k)$ term survives.  (This vanishing bracket is
the envelope mechanism that \cref{thm:envelope} elevates to a theorem: the
root sits exactly where the integrand of the combined asset-and-cash
ledger is zero, so moving it is free to first order.)

\eqref{eq:csecond}: differentiate \eqref{eq:cprime}, now keeping the root
term, with $\dd z_k/\dd k=1/x'(z_k)$ and $\Lambda'=\rho$:
\[
c''(k)=-e^{k}(1-u_k)+e^{k}\,\rho(z_k)\,\frac{\dd z_k}{\dd k}
=-e^{k}(1-u_k)+e^{k}\,\frac{\rho(z_k)}{x'(z_k)} .
\]

\eqref{eq:bl_division}: subtract \eqref{eq:cprime} from \eqref{eq:csecond}
and multiply by $e^{-k}$; the result $\rho(z_k)/x'(z_k)$ is $f_X(k)$ by
\eqref{eq:density}.  This is Breeden--Litzenberger \eqref{eq:BL}
transported to log-strike: $e^{-k}(c''-c')=f_X(k)$, equivalently
$\partial^2c/\partial y^2=f_Y(y)$ at $y=e^k$.
\end{proof}

Equation \eqref{eq:bl_division} deserves its plain-language reading: on
the grid, the risk-neutral density is \emph{one division} of two arrays
the build already holds, $\rho$ and $x'=e^{g\circ\Lambda}$.  No
differentiation of prices, no finite differences, no smoothing: the
density was never anything but these two arrays, and the call curve was
derived \emph{from} it, not the other way around.  This is the practical
content of choosing coordinates on laws rather than on curves.

Two further remarks.  First, \eqref{eq:cprime} identifies the
\emph{digital}: $-e^{-k}c'(k)=1-u_k=\Prob(Y>e^{k})$, so the model's
exercise probabilities are read off the same root.  Second, all five
identities hold for every admissible $\theta$ --- they are structure, not
calibration; the optimizer of \cref{sec:calibration} prices every trial
point through \eqref{eq:call} and nothing else.

\paragraph{A first ticket, on the solvable slice.}
The identities deserve one concrete run before they carry a real market.
Take the constant-speed slice of \cref{sec:constspeed}, with the scale
root-solved so that the ATM volatility is exactly $20\%$ at six months:
$s=\MacToyScale$, shift $m=\MacToyMu$ in closed form.  This is the slice
drawn in \cref{fig:transport}, and its two marked strike roots are now
priceable.  At the money, the forward's rank is
\MacToyAtmPercentile\% --- \emph{above} one half: the shift $m<0$ pushes
the median return below zero so that $e^{X}$ can average to one, which
is Jensen's inequality read off a percentile.  For the $k=0.10$ call:
the strike's rank is $u_k=\MacToyOtmPercentile\%$, the ledger holds
$G(z_k)=\MacToyShareOtm$ of the forward above that rank, the strike bill
on the same event is $e^{k}(1-u_k)=\MacToyCashOtm$, and
\eqref{eq:call} gives
\[
c(0.10)=\MacToyShareOtm-\MacToyCashOtm=\MacToyCallOtm
\qquad\Longrightarrow\qquad
\sigma_{\mathrm{imp}}=\MacToyIvOtmPct\%,
\]
about $0.6$ volatility points above the $20\%$ ATM level: the first
taste of a smile, produced by a symmetric law whose exponential tails
out-fatten the log-normal's.  \Cref{sec:examples} repeats this exercise
on a real listed option.

\subsection{The variance swap as the build's first moment}

The log contract gives the standard model-free variance proxy
\citep{Neuberger1994,DemeterfiDermanKamalZou1999,CarrMadan1998}: for a
continuous forward path monitored continuously, the fair strike of a
variance swap over $[0,T]$ is
\begin{equation}
w_{\mathrm{vs}}
=-2\,\E[\log Y]
=-2\,\E[X]
=-2\int_{-\infty}^{\infty}x(z)\,\rho(z)\,\dd z ,
\label{eq:varswap}
\end{equation}
one more integral over the arrays of the same build, with integrand
decaying like $|z|e^{-|z|}$.  Where the classical route replicates the log
contract from a strip of option prices across strikes
\citep{DemeterfiDermanKamalZou1999}, here the distribution is available
directly, so the strip replication is bypassed --- \eqref{eq:varswap} is
exact for the fitted law, and its parameter derivative rides the
Jacobian pass of \cref{sec:jacobian} at no extra build.

The honesty clause attaches immediately: \eqref{eq:varswap} prices the
\emph{log contract}.  Its identification with a realized-variance swap
assumes continuous monitoring of a continuous path; jumps and discrete
sampling take the usual corrections, which are outside the scope of a
single-expiry static model.  The paper therefore reports
$\sigma_{\mathrm{vs}}=\sqrt{w_{\mathrm{vs}}/\tau}$ as ``the var-swap
strike of the fitted law'' and claims nothing about realized variance.

\subsection{Five uses of one array}

It is worth pausing to collect the thread, because the pattern is the
paper's argument in miniature.  One quadrature build produces the arrays
$x(z_j)$, $x'(z_j)$, $G(z_j)$, $G'(z_j)$ on a fixed grid.  From these:
pricing is \eqref{eq:call} --- two lookups and a root; the density is
\eqref{eq:bl_division} --- one division; the variance swap is
\eqref{eq:varswap} --- one dot product; the calibration Jacobian is the
nodal parameter-sensitivity of $G$ evaluated at the same roots ---
\cref{thm:envelope}; and calendar order across expiries is the pointwise
comparison of two such ledgers --- \cref{thm:ledgerorder}.  Five
computations, one array.  A risk system built on this model holds one
object per slice and interrogates it; nothing is ever re-derived from
prices, because nothing was ever derived from prices in the first place.
